In the era of Industrial 4.0, smart manufacturing relies heavily on reliable fault diagnosis to ensure operational safety and system efficiency. However, most existing data-driven methods depend on the assumption that training and testing data share the same label space, which fails when encountering novel or unseen fault types in real-world industrial processes. Zero-shot learning (ZSL) has emerged as a promising paradigm to address this data scarcity by leveraging semantic information to bridge the gap between seen and unseen classes.

Despite their potential, current ZSL-based fault diagnosis methods often rely on hand-crafted binary attributes, which lack the granularity to capture complex fault dynamics, and frequently suffer from the domain shift problem between synthesized and real fault features. To address these challenges, we propose a Language-Driven Latent Diffusion for Continuous Semantic Alignment (Diff-LM-GZSL). By replacing discrete attributes with continuous semantic vectors extracted from pre-trained language models, our method provides a richer representation for fault descriptions and enhances the generalization capability of the diagnostic model.

The proposed Diff-LM-GZSL framework is evaluated on three benchmark industrial datasets: the Tennessee Eastman Process (TEP), a Hydraulic System dataset, and the Case Western Reserve University (CWRU) bearing dataset. Extensive experimental results demonstrate that our approach consistently outperforms nine state-of-the-art baseline models in generalized zero-shot fault diagnosis tasks. These findings validate the effectiveness of integrating language-driven latent diffusion with semantic alignment to achieve robust and high-fidelity unseen fault feature generation in complex industrial environments.